\documentclass[10pt]{amsart}

\usepackage[margin=1in]{geometry}
\usepackage[T1]{fontenc}
\usepackage{lmodern}
\usepackage{amsmath,amssymb,amsthm,mathtools}
\usepackage{microtype}
\usepackage{url}
\usepackage{enumitem}
\usepackage{alltt}
\usepackage{booktabs}
\usepackage[dvipsnames]{xcolor}
\usepackage[colorlinks=true,linkcolor=MidnightBlue,citecolor=OliveGreen,urlcolor=MidnightBlue]{hyperref}
\usepackage[capitalize,nameinlink]{cleveref}

\setlist{leftmargin=*,itemsep=2pt,topsep=4pt}
\emergencystretch=2.5em

\theoremstyle{plain}
\newtheorem{theorem}{Theorem}[section]
\newtheorem{lemma}[theorem]{Lemma}
\newtheorem{proposition}[theorem]{Proposition}
\newtheorem{corollary}[theorem]{Corollary}
\newtheorem{importedtheorem}[theorem]{Imported Theorem}
\theoremstyle{definition}
\newtheorem{definition}[theorem]{Definition}
\newtheorem{example}[theorem]{Example}
\theoremstyle{remark}
\newtheorem{remark}[theorem]{Remark}
\crefname{importedtheorem}{Imported Theorem}{Imported Theorems}
\numberwithin{equation}{section}

\newcommand{\A}{\textup{[A]}}
\newcommand{\Oa}{\textup{[O]}}
\newcommand{\N}{\mathbb{N}}
\newcommand{\Z}{\mathbb{Z}}
\newcommand{\tl}{\tau^{<}}
\newcommand{\rate}{2^{O(n/\sqrt{\log n})}}
% Lean names: typewriter, with line breaks allowed after '.' and '_'.
\DeclareUrlCommand\lean{\urlstyle{tt}\def\UrlBreaks{\do\.\do\_}\def\UrlBigBreaks{}}
% Repository names: typewriter, with line breaks allowed after '/' and '-'.
\DeclareUrlCommand\repo{\urlstyle{tt}\def\UrlBreaks{\do\/\do\-}\def\UrlBigBreaks{}}

\begin{document}

\title[The parity of $\pi(x)$ in time $x^{o(1)}$]{The parity of the prime-counting function in time $x^{o(1)}$,\\ assuming the new $\#$SAT algorithms}

\author{Vaibhav Tulsyan}

\date{November 2026}

\subjclass[2020]{11Y11, 11Y16, 11N05, 68Q25, 03B35}
\keywords{Prime-counting function, parity, squarefree numbers, M\"obius function, deterministic prime finding, exact counting, Exponential Time Hypothesis, Lean~4, formal verification}

\begin{abstract}
Computing the prime-counting function $\pi(x)$ exactly takes $x^{1/2+o(1)}$ steps by the analytic
method of Lagarias and Odlyzko, and to our knowledge its parity is known faster only on short intervals: Tao,
Croot and Helfgott decide whether an interval of length at most $x^{1/2+c}$ near $x$ contains an odd
number of primes in time $x^{1/2-c+o(1)}$. However, their identity needs the M\"obius function
$\mu(d)$ for all $d\le\sqrt x$, and we know of no bound of the form $x^{o(1)}$ for the parity, even
conditionally. We present deterministic algorithms that output
$\pi(x)\bmod2$ and the parity of the number of primes in any interval $[A,B]$ in time
$2^{O(n/\sqrt{\log n})}=x^{o(1)}$, where $n$ is the input bit length. Our approach replaces $\mu(d)$ by
its residue modulo $2$, which is the squarefree indicator again: the number of chains
$d_1^2\mid m$, $d_2^2\mid d_1,\dots$ of depth $L$ is odd exactly when $m$ is squarefree, for
$m<2^{2^L}$. The parity of $\pi$ on $[A,B]$ is then the parity of one count of tuples satisfying
$O(\log n)$ multiplications, which Anthropic's counting engine evaluates through a network whose
variables have geometrically shrinking widths. We prove both theorems in Lean~4
(\texttt{piParity\_TM\_rate}, \texttt{piIntervalParity\_TM\_rate}; axioms \texttt{propext},
\texttt{Classical.choice}, \texttt{Quot.sound}). The result does not give deterministic prime finding
or the exact value of $\pi(x)$: the parity can be even on every interval one queries, and exact counts
seem to need iterated products, for whose product-tree verifiers we prove a lower bound in Lean. The
constants are astronomical, and the result rests on Anthropic's Lean development of its claimed
refutation of the Exponential Time Hypothesis, which is under review.
\end{abstract}

\maketitle

% =====================================================================
\section{Introduction}\label{sec:intro}
% =====================================================================

How fast can one count primes? The prime-counting function $\pi(x)$ is computed exactly in
$x^{2/3+o(1)}$ steps by the combinatorial method of Meissel, Lehmer, Lagarias, Miller and
Odlyzko~\cite{LMO85,DR96}, and in $x^{1/2+o(1)}$ steps by the analytic method of Lagarias and
Odlyzko~\cite{LO87}. The same bound $x^{1/2+o(1)}$ is the best known for deterministically
\emph{finding} a prime between $x$ and $2x$, the question of the Polymath4 project~\cite{Tao09,TCH12}.
Tao, Croot and Helfgott~\cite{TCH12} showed that one bit of $\pi$ can be had faster: they decide
whether an interval $[a,b]\subseteq[N,2N]$ of length at most $N^{1/2+c}$ contains an odd number of
primes in time $O(N^{1/2-c+o(1)})$, for an absolute $c>0$ (their Theorem~1.2), and
Lelechenko~\cite{Lel13} made $c$ explicit.

However, both results are for short intervals, and the obstacle to longer ones is structural. The
identity behind~\cite{TCH12} sums a divisor-sum function over $d\le\sqrt N$ with weights $\mu(d)$, and
$\mu(d)$ for every $d\le\sqrt N$ is itself an $N^{1/2}$ computation. We are not aware of a published bound
below $x^{1/2+o(1)}$ for $\pi(x)\bmod 2$; Sladkey~\cite[ref.~5]{Sla12} cites a manuscript by Oliveira e
Silva on computing the parity of $\pi$ as in preparation, and we could not find it.

In October 2026 OpenAI~\cite{OpenAI26} and Anthropic~\cite{Anthropic26} (below \Oa{} and \A) each
claimed a deterministic algorithm that counts the satisfying assignments of a $k$-CNF in subexponential
time, refuting the Exponential Time Hypothesis; both claims come with Lean~4 developments and are under
review. Rudich and Rousseau~\cite{RR26} used the engines to factor in $N^{o(1)}$, and our earlier
notes~\cite{TulFac26,TulCN26} proved in Lean, on \A's engine, bounds of the form $\rate$ for factoring,
class numbers, square roots and divisor sums. We found no derivation of a parity result for $\pi$ from
either refutation (\cref{sec:related}).

\subsection{The problem and the running example}
We consider the \emph{interval prime parity} problem: given $A$ and $B$, output
$\pi([A,B])\bmod 2$, where $\pi([A,B])=\#\{p\text{ prime}:A\le p\le B\}$. The case $A=1$, $B=x$ is
$\pi(x)\bmod2$. We write $\tl(m)=\#\{(a,b):ab=m,\ a<b\}$ and
\[
  C_L(m)=\#\{(d_1,\dots,d_L):\ d_1^2\mid m,\ d_2^2\mid d_1,\ \dots,\ d_L^2\mid d_{L-1}\}.
\]

\begin{example}[Running example]\label{ex:ten}
Take $[A,B]=[1,10]$ and $L=2$. \Cref{tab:ten} lists $\tl(m)$, $C_2(m)$ and their product. The sum of
the last row is $16$, which is even, and indeed $\pi(10)=4$. The summand is odd exactly at the primes
$2,3,5,7$: at $m=6$ and $m=10$ the factor $\tl(m)=2$ is even, and at $m=4,8,9$ the factor $C_2(m)=2$
is even.
\end{example}

\begin{table}[ht]
\centering\small
\begin{tabular}{@{}lcccccccccc|c@{}}
\toprule
$m$ & 1 & 2 & 3 & 4 & 5 & 6 & 7 & 8 & 9 & 10 & sum\\
\midrule
$\tl(m)$ & 0 & 1 & 1 & 1 & 1 & 2 & 1 & 2 & 1 & 2 & \\
$C_2(m)$ & 1 & 1 & 1 & 2 & 1 & 1 & 1 & 2 & 2 & 1 & \\
$\tl(m)\,C_2(m)$ & 0 & 1 & 1 & 2 & 1 & 2 & 1 & 4 & 2 & 2 & 16\\
\midrule
$m$ prime & & $\ast$ & $\ast$ & & $\ast$ & & $\ast$ & & & & 4\\
\bottomrule
\end{tabular}
\caption{The running example: the one-count identity on $[1,10]$ with $L=2$.}\label{tab:ten}
\end{table}

\subsection{The naive route and the principle}
At first one might count primes exactly, through a signed identity such as
$[m\text{ squarefree}]=\sum_{d^2\mid m}\mu(d)$, and read off the parity. However, a counting engine
sums non-negative table entries over labellings; a weight $\mu(d)$ on a label $d$ of $\Theta(n)$ bits
needs the factorisation of $d$ inside one node, and the nodes of \A's networks hold $o(n)$ bits. We
call this the \emph{M\"obius barrier}. Instead, we rely on a different principle: \emph{modulo $2$,
$\mu(d)\equiv[d\text{ squarefree}]$, so the M\"obius weight is the problem itself at half the size, and
the recursion closes after $\log_2 n$ halvings.} The resulting chain count is a plain count of
divisibility witnesses, and the witnesses shrink geometrically, so their total size is $O(n)$.

\subsection{What we prove}
In this paper we prove that one count of a fixed system of $O(\log n)$ multiplications has the parity of
$\pi([A,B])$, build the network that evaluates the count with $O(n)$ shared bits, and prove in Lean that
one Turing machine per problem outputs the parity.

\begin{theorem}[\cref{sec:statement}; machine-checked]\label{thm:main}
There are Turing machines $M_1,M_2$ in Mathlib's \lean{TM2} model, with finite alphabets on all stacks,
and for each of them a constant $C\in\N$, such that on every input word $w$ of length $n$ the machine halts within
$\lfloor C\cdot2^{Cn/\sqrt{\log_2(n+2)}}\rfloor$ steps with the following output:
\begin{enumerate}[label=(\roman*)]
\item on $w=$ \lean{encodeNat}~$x$, the machine $M_1$ outputs \lean{encodeNat}~$(\pi(x)\bmod2)$
(\lean{piParity_TM_rate});
\item on $w=$ \lean{encodePair}~$A$~$B$, the machine $M_2$ outputs
\lean{encodeNat}~$(\pi([A,B])\bmod2)$ (\lean{piIntervalParity_TM_rate}).
\end{enumerate}
\end{theorem}

\noindent Since $n=\log_2x+O(1)$ in (i), the bound is $x^{O(1/\sqrt{\log\log x})}=x^{o(1)}$, against
$x^{1/2+o(1)}$ for the exact $\pi(x)$~\cite{LO87}. In (ii) it is $\max(A,B)^{o(1)}$ for every interval,
against $N^{1/2-c+o(1)}$ for intervals of length at most $N^{1/2+c}$ in $[N,2N]$~\cite{TCH12}. To our
knowledge, no bound of the form $x^{o(1)}$ was known for either problem, even conditionally.

\subsection{Contributions}
Our contributions are:
\begin{itemize}
\item We prove that the chain count $C_L(m)$ is odd exactly when $m$ is squarefree, for $1\le m<2^{2^L}$,
and that $\sum_{m\in[A,B]}\tl(m)C_L(m)\equiv\pi([A,B])\pmod2$ for $B<2^{2^L}$
(\cref{sec:chains}).
\item We present a leveled description of this count, with $3L+3$ variables at levels $0,\dots,L$, and a
version of our network library in which a level-$k$ variable has about $n/2^k$ bits, so that
$\Theta(\log n)$ multiplications share $O(n)$ bits in total (\cref{sec:network}).
\item We prove the two theorems of \cref{thm:main} in Lean and list what a reader must trust
(\cref{sec:statement,sec:formal}).
\item We show what the result does not give: an adversary argument (on paper) for why parities alone do
not find primes, and a Lean-checked lower bound for product-tree verifiers of the iterated product that
exact counting by this route would need (\cref{sec:not}).
\end{itemize}
The number-theoretic content is small: the chain lemma, and the observation that working modulo $2$
bypasses the M\"obius barrier. The network is engineering on top of our earlier library~\cite{TulCN26};
it is what makes the rate hold.

\subsection{Status of each claim}
We separate four kinds of claim. \emph{Machine-checked in Lean}: the two theorems of \cref{thm:main},
with the axioms \lean{propext}, \lean{Classical.choice} and \lean{Quot.sound} only, and every lemma whose
Lean name we give, including the tree bound of \cref{thm:tree}. \emph{Proved on paper only}: the
adversary statement of \cref{prop:adv}; the link from network cuts to the tree model, and the circuit
simulation, in \cref{sec:barrier}, where the need for long products is argued informally; the modulo-$k$ variant of \cref{rem:modk}; and the comparisons with earlier bounds.
\emph{Assumed}: nothing beyond the three axioms; \A's results enter as Lean theorems of its development,
which the kernel checks with ours. \emph{Checked numerically}: \cref{ex:ten}; the chain lemma for all
$1\le m<\min(2^{2^L},5000)$, $L\le3$, and the one-count identity on every interval $[A,B]$ with
$1\le A\le300$, $A\le B\le3000$, $L=4$, by a short exact computation that is not part of the proof; and
the same identities for $m\le255$ and intervals in $[1,120]$ by \lean{#eval} in Lean.

\subsection*{Outline}
\Cref{sec:engine} recalls \A's counting theorem. \Cref{sec:chains} proves the chain lemma and the
one-count identity, \cref{sec:network} builds the network and the machine, and \cref{sec:statement}
quotes the statements. \Cref{sec:not} says what the result does not give,
\cref{sec:formal} describes the Lean development, \cref{sec:related,sec:limitations} cover related
work and limitations, and \cref{sec:conclusion} concludes.

% =====================================================================
\section{The engine counts accepted labellings}\label{sec:engine}
% =====================================================================

In this section we recall the one theorem of \A{} that we compose with, in the form used by our
factoring note~\cite[\S2]{TulFac26}. We use \A's Lean development at commit \texttt{16a5ce4}
(\texttt{anthropics/formal-math}, folder \texttt{eth}, Lean v4.33.1, Mathlib \texttt{0df444a}).

A \emph{network on $[g]$ with parameters $P,B,\tau$} has, for every node $a$, settings of at most $P$
bits and a weight $W_a$, and for every pair of nodes an interaction that factors through $c_{a,b}$
shared bits, with $\sum_bc_{a,b}\le P$; its value is $Z=\sum_z\prod_aW_a(z_a)\prod_{\{a,b\}}\Psi_{a,b}$ and
its total is $U=\sum_{a\ne b}c_{a,b}$. We use only \emph{label systems}: wires placed on pairs of
nodes, each carrying a label seen by both ends, with $W_a$ a $0/1$ table. The value is then the
number of labellings accepted by every table.

\begin{importedtheorem}[\A, Lemma~8(i) and Theorem~A(b); see~{\cite[Imported Thm.~2.2]{TulFac26}}]\label{imp:engine}
Let $K=K_\ell\ge2$ be a level of \A's schedule, $g=3K^8$, and $e_\ell\le(3x_0+6)/\sqrt{\log_2K}$ for an
absolute $x_0$. A network on $[g]$ with parameters $P,B,\tau$ and total $U$ can be evaluated within
$\bigl(K^g+2^{e_\ell K^2L}(K(L+B'))^{C_0K^2}2^{2gP/K}\tau\bigr)(g(P+B))^{C_1}$ operations, where
$L=\lceil(U+gP/K)/K^2\rceil$, $B'=\lceil2gP/K\rceil+g^2B$ and $C_0,C_1$ are absolute.
\end{importedtheorem}

\noindent So a network family with $U=O(n)$ and $gP=O(n)$ is evaluated in $\rate$ steps at the level
$K^{40}\lesssim n$ of~\cite[\S5]{TulFac26}. The cost is exponential in $U/\sqrt{\log n}$: a
network with $U=\Theta(n\log n)$ would cost $2^{\Theta(n\sqrt{\log n})}$, which is superpolynomial in $x$.
Everything below is about keeping $U$ linear.

% =====================================================================
\section{Squarefree chains and the one-count identity}\label{sec:chains}
% =====================================================================

In this section we prove the two counting facts. Both are in the pure Mathlib module
\lean{ChainParity}, whose definitions are
\begin{alltt}\footnotesize
def chainCount : \(\mathbb{N}\) \(\to\) \(\mathbb{N}\) \(\to\) \(\mathbb{N}\)
  | 0, _ => 1
  | L+1, n => \(\sum\) d \(\in\) n.divisors.filter (fun d => d * d \(\mid\) n), chainCount L d
def tauLt (m : \(\mathbb{N}\)) : \(\mathbb{N}\) := ((Nat.divisorsAntidiagonal m).filter fun p => p.1 < p.2).card
\end{alltt}
so \lean{chainCount}~$L$~$n=C_L(n)$ for $n\ge1$ and \lean{tauLt}~$m=\tl(m)$.

\begin{lemma}[Chain lemma; \lean{chainCount_mod_two}]\label{lem:chain}
For $L\ge0$ and $1\le n<2^{2^L}$, $C_L(n)\equiv[n\text{ squarefree}]\pmod2$. In Lean:
\begin{alltt}\footnotesize
theorem chainCount_mod_two \{L n : \(\mathbb{N}\)\} (h1 : 1 \(\le\) n) (h : n < 2 ^ 2 ^ L) :
    chainCount L n % 2 = if Squarefree n then 1 else 0
\end{alltt}
\end{lemma}

\begin{proof}
Write $n=st^2$ with $s$ squarefree (\lean{Nat.sq_mul_squarefree_of_pos}). Then $d^2\mid n$ if and only
if $d\mid t$ (\lean{dvd_of_mul_self_dvd}), and $n$ is squarefree if and only if $t=1$. Read the identity
$\sum_{d\mid t}\mu(d)=[t=1]$ modulo $2$, where $\mu(d)\equiv[d\text{ squarefree}]$:
\begin{equation}\label{eq:mu2}
  \#\{d:\ d^2\mid n,\ d\text{ squarefree}\}\ \equiv\ \sum_{d\mid t}\mu(d)\ =\ [n\text{ squarefree}]\pmod 2
\end{equation}
(\lean{key_zmod}). Now induct on $L$. For $L=0$, $n<2$ forces $n=1$, and $C_0(1)=1$. For $L+1$,
$C_{L+1}(n)=\sum_{d^2\mid n}C_L(d)$, and every such $d$ satisfies $d\le\sqrt n<2^{2^L}$, so by induction
$C_L(d)\equiv[d\text{ squarefree}]$, and~\eqref{eq:mu2} finishes.
\end{proof}

\noindent The hypothesis $n<2^{2^L}$ cannot be dropped: $C_1(16)=3$ counts $d_1\in\{1,2,4\}$, but $16$
is not squarefree, while $C_2(16)=C_1(1)+C_1(2)+C_1(4)=4$ is even. The depth must grow like $\log_2 n$.

\begin{proposition}[One-count identity; \lean{primeParity}]\label{prop:onecount}
If $B<2^{2^L}$, then $\sum_{m=A}^{B}\tl(m)\,C_L(m)\equiv\pi([A,B])\pmod2$. In Lean:
\begin{alltt}\footnotesize
theorem primeParity \{L A B : \(\mathbb{N}\)\} (hB : B < 2 ^ 2 ^ L) :
    (\(\sum\) m \(\in\) Finset.Icc A B, tauLt m * chainCount L m) % 2 =
      ((Finset.Icc A B).filter Nat.Prime).card % 2
theorem primeCounting_eq_card (x : \(\mathbb{N}\)) :
    Nat.primeCounting x = ((Finset.Icc 1 x).filter Nat.Prime).card
\end{alltt}
\end{proposition}

\begin{proof}
We show that each summand is odd exactly when $m$ is prime (\lean{term_zmod}). For $m=0$ and $m=1$,
$\tl(m)=0$. For $m\ge2$ not squarefree, $C_L(m)$ is even by \cref{lem:chain}. For $m\ge2$ squarefree,
$C_L(m)$ is odd, and the swap $(a,b)\mapsto(b,a)$ pairs the factorisations with $a<b$ and those with
$a>b$, with none on the diagonal, so $\tau(m)=2\tl(m)$ (\lean{card_divisors_eq_two_tauLt}). Hence
$\tl(m)=2^{\omega(m)-1}$, which is odd if and only if $\omega(m)=1$, that is, $m$ is prime
(\lean{tauLt_prime}, \lean{tauLt_even}). The second statement relates Mathlib's \lean{Nat.primeCounting}
to the count on $[1,x]$.
\end{proof}

\noindent In \cref{ex:ten}, $L=2$ suffices since $10<2^{2^2}$. The machine of \cref{sec:network} takes
$L$ to be the bit length of $n$, here $L=3$; on $[1,10]$ the values $C_3(m)$ equal $C_2(m)$, and the
count is again $16$.

% =====================================================================
\section{One count of a leveled network}\label{sec:network}
% =====================================================================

In this section we write the sum of \cref{prop:onecount} as the number of solutions of a system of
products, explain why our earlier network library cannot evaluate it at the rate, describe the leveled
version that can, and assemble the machine.

\subsection{The description}
A \emph{description}, in the sense of our class-number note~\cite[\S3]{TulCN26}, lists variables,
products of a variable by a variable or a constant, chains (equations between sums of products) and
comparisons with constants or other variables. For depth $L$ and interval $[A,B]$ the description
\lean{primeDescL} has, for $k=0,\dots,L$, the variables $d_k,s_k,e_k$, with $d_0=m$, $s_0=a$, $e_0=b$,
and the constraints
\[
  ab=m,\qquad d_k\cdot d_k=s_k,\quad s_k\cdot e_k=d_{k-1}\ \ (1\le k\le L),\qquad
  m\ge\max(A,1),\quad m\le B,\quad a<b .
\]
So $s_k=d_k^2$ and $e_k=d_{k-1}/d_k^2$ witness $d_k^2\mid d_{k-1}$, and they are determined by the
chain. A solution is a factorisation $m=ab$ with $a<b$ together with a chain of depth $L$ from $m$, so
the number of solutions is $\sum_{m\in[A,B]}\tl(m)C_L(m)$ (\lean{primeDescL_count}). The condition
$m\ge1$ excludes $m=0$, which every $d$ would divide. In \cref{ex:ten} the $16$ solutions include, for
$m=8$, the four tuples with $(a,b)\in\{(1,8),(2,4)\}$ and $d_1\in\{1,2\}$, $d_2=1$.

\subsection{Why the library of the class-number note does not suffice}
ArithNet~\cite{TulCN26} gives every variable the full box $[0,2^{n_B})$ with $n_B>n$ bits, and its
parameter bound is $U\le2^{11}c_A(D)(n+2)$ with $c_A(D)=64(\mathrm{nV}+\mathrm{nP}+\mathrm{nQ}+1)(\mathrm{nC}+1)$.
Our description has $\mathrm{nV}=3L+3$ variables, $\mathrm{nP}=4L+2$ products and $\mathrm{nC}=2L+1$
chains, so $c_A(D)=\Theta(L^2)$ and the bound is $U=\Theta(n\log^2n)$, which makes the engine cost
$2^{\Theta(n\log^{3/2}n)}$, worse than trial division. Even with every per-chain copy removed, $3L+3$
variables of $n$ bits carry $\Theta(n\log n)$ bits, which costs $2^{\Theta(n\sqrt{\log n})}$. The
witnesses themselves are small: along a chain, $d_k^{2^k}\le m<2^n$ (\lean{chain_pow_le}), so $d_k$ has at
most $n/2^k$ bits and all of them together fewer than $2n$. The network should pay for that.

\subsection{Levels}
A \emph{leveled description} (\lean{LDesc}, module \lean{ArithLDesc}) is a description together with a
level for every variable and every chain. A variable of level $k$ ranges over $[0,2^{b\ell_k})$, where
$b$ is the limb width of ArithNet and $\ell_k=\ell/4^{\min(\lfloor k/2\rfloor,\,r-1)}$ is the number of
limbs at level $k$ (here $\ell=4^{r-1}$ is the number of limbs of a full-width variable in~\cite{TulFac26}); \lean{lbox_spec} shows that a
level-$k$ variable $y$ with $y^{2^k}<2^n$ fits. A level-$k$ chain is checked by the two-dimensional
Sch\"onhage--Strassen pipeline of~\cite{TulFac26} at the transform shape of level $k$, with the limb width
and the Fermat modulus of the top level, so a level-$k$ variable is a prefix of the top-level limb
format and the levels need no glue. Well-formedness (\lean{LDesc.WFB}) asks that the operands of a
product have level at least that of its chain, that comparisons be between variables of one level or
against constants, and that variables enter chains only as products with the constant~$1$. In
\lean{primeDescL}, $m,a,b$ have level $0$, $d_k$ has level $k$, $s_k,e_k$ have level $k-1$, and the chains
$d_k\cdot d_k=s_k$ and $s_k\cdot e_k=d_{k-1}$ have level $k-1$.

Every object (variable, product, chain, comparison) has its own family of wires and vertices, so the
cost is a sum over objects of a function of the level. The number of limbs is geometric,
$\sum_{k<K}\ell_k\le4\ell+K$ (\lean{sum_kl_le}), and a description whose $i$-th object of each kind has
level at least $i/c$ (\lean{LDesc.Spread}~$c$; \lean{primeDescL_spread} gives $c=8$) has $O(c)$ objects
per level.

\begin{proposition}[\lean{arithNet_value}, \lean{arithSys_U_le}, \lean{arithSys_is}]\label{prop:arithl}
Let $g\ge2$ and let $D$ be a leveled description that is well formed at size $n$, has spread $c\ge1$ and
at most $\lfloor\log_2(n+2)\rfloor+3$ levels. The network of $D$ has value equal to the number of solutions
of $D$ in its boxes, total $U\le2^{19}c\,(n+2)$, and node load
$P=2^{22}c\bigl((n+2)/g+d(n)\bigr)$ with $d(n)=O(n^{3/4})$ the vertex degree of~\cite{TulCN26}.
\end{proposition}

\noindent For \lean{primeDescL} at depth $L=\lfloor\log_2n\rfloor+1$ this gives $U\le2^{22}(n+2)$. We call
this version of the library ArithL. It is a re-parameterisation of ArithNet with one more coordinate,
not a new idea; its proofs reuse the factoring and class-number pipelines unchanged, and its value lies
in the bookkeeping that keeps $U$ linear for $\Theta(\log n)$ multiplications.

\subsection{The machine}
A counting package~\cite[\S4]{TulCN26} provides, for every input word $w$ and query word $\pi$, a network
whose value is $Q(w,\pi)$, with $U$ and $gP$ linear in $|w|$ and polynomial-time readers. From any
polynomial-time map $w\mapsto D(w)$ to well-formed leveled descriptions with spread $c$,
\lean{ArithLSpec.pkg_exists} builds a package with constant $2^{22}c$ and
$Q(w,\pi)=\mathrm{count}(D(w))$. We use it twice, with $n=|w|$: \lean{arithPi} reads $x$ from $w$ and takes
$[A,B]=[1,x]$, and \lean{arithPiInt} reads $A$ and $B$ from the two halves of \lean{encodePair}~$A$~$B$.
Both take $L$ to be the bit length of $n$, so that $B<2^n<2^{2^L}$ (\lean{primeL_spec}), and
\lean{primeDesc_count_parity} gives
$\mathrm{count}\bmod2=\pi([A,B])\bmod2$.

The query specification \lean{specPi} asks the single query $\pi=[\,]$ and outputs the binary word of
the count modulo $2$. The generic one-query machine \lean{GR.queryM_TM} of~\cite[Thm.~4.1]{TulCN26}
turns a package and a polynomial-time specification with output of length $O(|w|)$ into a finite-alphabet
Turing machine that outputs the specified answer within $\lfloor C\cdot2^{C|w|/\sqrt{\log_2(|w|+2)}}\rfloor$
steps. Composing the two gives \lean{piParity_TM_of} and \lean{piIntervalParity_TM_of}, which the module
\lean{Final} restates as the theorems of \cref{thm:main}.

% =====================================================================
\section{What a reader must trust}\label{sec:statement}
% =====================================================================

A reader who wants to trust \cref{thm:main} must check the following definitions and nothing else of
ours. The statements, verbatim from \lean{Final} up to line breaks, are below; we render Unicode
symbols as mathematics.
\begin{alltt}\footnotesize
theorem piParity_TM_rate :
    \(\exists\) (M : Turing.TM2ComputableAux Bool Bool) (C : \(\mathbb{N}\)), (\(\forall\) i, Finite (M.tm.\(\Gamma\) i)) \(\wedge\)
      \(\forall\) x : \(\mathbb{N}\), GraphProblems.Outputs M (Computability.encodeNat x)
        (Computability.encodeNat (Nat.primeCounting x % 2))
        \(\lfloor\)(C : \(\mathbb{R}\)) * 2 ^ ((C : \(\mathbb{R}\)) * (Computability.encodeNat x).length /
          Real.sqrt (Real.logb 2 ((Computability.encodeNat x).length + 2)))\(\rfloor\sb{+}\)
theorem piIntervalParity_TM_rate :
    \(\exists\) (M : Turing.TM2ComputableAux Bool Bool) (C : \(\mathbb{N}\)), (\(\forall\) i, Finite (M.tm.\(\Gamma\) i)) \(\wedge\)
      \(\forall\) A B : \(\mathbb{N}\), GraphProblems.Outputs M (encodePair A B)
        (Computability.encodeNat (((Finset.Icc A B).filter Nat.Prime).card % 2))
        \(\lfloor\)(C : \(\mathbb{R}\)) * 2 ^ ((C : \(\mathbb{R}\)) * (encodePair A B).length /
          Real.sqrt (Real.logb 2 ((encodePair A B).length + 2)))\(\rfloor\sb{+}\)
\end{alltt}
The only definitions of ours they use are those of the pair encoding (module \lean{ModP.Encode}
of~\cite{TulCN26}):
\begin{alltt}\footnotesize
def padR (n : \(\mathbb{N}\)) (l : List Bool) : List Bool := l ++ List.replicate (n - l.length) false
def encodePair (m a : \(\mathbb{N}\)) : List Bool :=
  padR (max (encodeNat m).length (encodeNat a).length) (encodeNat m) ++
    padR (max (encodeNat m).length (encodeNat a).length) (encodeNat a)
\end{alltt}
\texttt{Outputs} is the one-line definition of \A's statement file, which we also trust:
\begin{alltt}\footnotesize
def Outputs (M : TM2ComputableAux Bool Bool) (x y : List Bool) (T : \(\mathbb{N}\)) : Prop :=
  Nonempty (TM2OutputsInTime M.tm (x.map M.inputAlphabet.invFun)
    (some (y.map M.outputAlphabet.invFun)) T)
\end{alltt}
Behind these are Mathlib's \lean{TM2OutputsInTime}, \lean{Computability.encodeNat} (little-endian
binary, the empty word for $0$), \lean{Nat.primeCounting}, \lean{Nat.Prime}, \lean{Finset.Icc},
\lean{Real.logb} and \lean{Real.sqrt}, from Mathlib \texttt{0df444a}. The definitions \lean{chainCount}
and \lean{tauLt} of \cref{sec:chains} do not occur in the statements; they are part of the proof.

% =====================================================================
\section{What the parity does not give}\label{sec:not}
% =====================================================================

In this section we say plainly what does not follow: a deterministic prime finder, the exact value of
$\pi(x)$, and anything practical (the last is in \cref{sec:limitations}).

\subsection{Prime finding}
Bisection finds a prime from any interval with an odd number of primes: one of its halves again has an
odd number, and after $O(n)$ queries the interval is a single prime. So a fast parity oracle gives a fast
prime finder as soon as one knows \emph{some} interval of odd parity. Tao, Croot and Helfgott make the
same remark and add that they ``do not know of a method to quickly and deterministically locate an
interval with an odd number of primes''~\cite[\S1]{TCH12}. Our oracle is faster, but it does not change
this. The parity of $\pi$ on $[x,2x]$ may be even, and nothing we know about the primes forces an odd
count on any explicit family of intervals of small size. The following statement makes the gap precise in
an oracle model; it is proved on paper only.

\begin{proposition}[Adversary; on paper]\label{prop:adv}
Let $H\ge4$. Consider deterministic algorithms that learn about an unknown set $S\subseteq[x,2x]$ only
by adaptive queries of two kinds, the parity of $|S\cap I|$ for an interval $I$, and membership $y\in S$,
and that are promised only that $S$ meets every interval of $H$ consecutive integers in $[x,2x]$. Every
such algorithm that outputs an element of $S$ makes at least $(H-3)/2$ queries on some admissible $S$.
\end{proposition}

\begin{proof}
Answer ``even'' to every parity query and ``no'' to every membership query, and suppose the algorithm
stops after $q<(H-3)/2$ queries, $q_{\mathrm p}$ of parity and $q_{\mathrm m}$ of membership, with
output $y$. The endpoints of the queried intervals cut $[x,2x]$ into at most $2q_{\mathrm p}+1$ cells, and every queried interval meets $[x,2x]$ in a union of cells. Let $S$ consist of
all integers of $[x,2x]$ except the membership-queried points and $y$; then, in each cell where $S$ has
an odd number of points, remove one more point of $S$. Every queried interval now meets $S$ in an even
number of points, every membership answer is correct, and $y\notin S$. At most
$q_{\mathrm m}+1+(2q_{\mathrm p}+1)\le2q+2<H$ integers of $[x,2x]$ are missing from $S$, so $S$ meets every
window of $H$ consecutive integers. So $S$ is admissible, consistent with all answers, and the output is
wrong.
\end{proof}

\noindent With $H=x^{0.525}$, the gap bound of Baker, Harman and Pintz~\cite{BHP01}, the matching
strategy (test the integers from $x$ upwards by a polynomial-time primality test~\cite{AKS04}) takes
$x^{0.525+o(1)}$ steps, worse than the $x^{1/2+o(1)}$ of~\cite{LO87}. The proposition does not say that
parities cannot find primes; it says that a proof that they do must use more about the primes than a
gap bound, and we know of no such argument. For residues modulo $M$ (\cref{rem:modk}) the same
argument, removing up to $M-1$ points per cell, gives about $H/(2M)$ queries.

\subsection{Exact counts, and the iterated-product barrier}\label{sec:barrier}
The exact $\pi([A,B])$ would give prime finding at once. By this route it needs an identity that isolates
$\omega(m)=1$ exactly, not modulo a small number. We argue informally, as in our working
notes~\cite[\texttt{research/prime-finding.md}, \S4]{TulsyanRepo26}: on squarefree $m$, the weights
built from $O(1)$ multiplications of divisors are combinations of $t^{\omega(m)}$ for small $t$; the
indicator of $\omega=1$ on $0\le\omega\le n/\log n$ needs $\Omega(n/\log n)$ such terms, unless one works
modulo $t^2$, which yields only residues. So exact counting asks for the $k$-fold divisor function with
$k=\Theta(n/\log n)$: a verifier of $m=d_1\cdots d_k$ with $\sum_i\operatorname{bits}(d_i)=n$. The product
tree verifies it with $U=\Theta(n\log k)=\Theta(n\log n)$, and then \cref{imp:engine} gives
$2^{\Theta(n\sqrt{\log n})}$, which is superpolynomial in $x$. A verifier with $U=o(n\sqrt{\log n})$ would
give the exact $\pi(x)$ in $x^{o(1)}$; we do not know one.

In the restricted model of product trees we can prove that no tree shape helps. A \emph{tree verifier}
multiplies the factors along a binary tree $T$ with leaf sizes $\ell_i$; its cost is
$\mathrm{cost}(T)=\sum_{v}\mathrm{wt}(v)=\sum_i\ell_i\,\mathrm{depth}(i)$, where $\mathrm{wt}(v)$ is the
total size of the leaves below the internal node $v$.

\begin{theorem}[\lean{BTree.cost_ge_entropy}, \lean{BTree.cost_ge_equal}; machine-checked]\label{thm:tree}
For every binary tree with positive leaf weights $\ell_i$ and $L=\sum_i\ell_i$,
$\mathrm{cost}(T)\ge\sum_i\ell_i\log_2(L/\ell_i)$. If all $k$ leaves have weight $w$, then
$\mathrm{cost}(T)\ge kw\log_2k=n\log_2k$ for $n=kw$.
\end{theorem}

\noindent The proof is an induction whose step is the bound $1$ on binary entropy
(\lean{Real.binEntropy_le_log_two}). The link to networks is on paper: in a tree verifier the subtree of
$v$ is separated from the rest by a cut across which the product indicator has rank $2^{\mathrm{wt}(v)}$
(\lean{rank_prodMat_pow}, Lean-checked), so by the cut-rank lemma
(\lean{rank_le_two_pow_of_eq_sum_mul}, module \lean{Barriers.RankBounds} of~\cite{TulsyanRepo26}) the cut carries at least $\mathrm{wt}(v)$ bits, and summing over $v$ gives
$U\ge\mathrm{cost}(T)$. A lower bound for \emph{all} networks is out of reach: the engine simulates any
Boolean circuit of size $O(n^{1+1/40})$ with $U$ of the order of its size, so $U=\omega(n)$ for every
verifier of $m=\prod_id_i$ would be a superlinear circuit lower bound for an explicit function, which is
open.

\begin{remark}[Residues modulo $k$; on paper]\label{rem:modk}
Modulo $k$, $\mu(d)\equiv\tau_{k-1}(d)[d\text{ squarefree}]$, since $\tau_{k-1}(d)=(k-1)^{\omega(d)}$ on
squarefree $d$. Iterating as in \cref{lem:chain} gives a $\tau_{k-1}$-weighted chain $W_k(m)\equiv
[m\text{ squarefree}]\pmod k$, and with $W=W_{k^2}$ one gets
$\sum_{m\in I}\tau_k(m)W(m)\equiv[1\in I]+k\,\pi(I)\pmod{k^2}$, hence $\pi(I)\bmod k$. We checked this
by \lean{#eval} for $k=2,3$ on all intervals in $[1,100]$. The verifier needs $(k^2-1)$-fold products,
so its cut grows like $n\log k$; we have not built it and claim no running time. Residues modulo
small $M$ do not escape \cref{prop:adv}.
\end{remark}

% =====================================================================
\section{The Lean development}\label{sec:formal}
% =====================================================================

In this section we list the modules and state the trusted base. The development is in the public
repository~\cite{TulsyanRepo26}: directories \texttt{lean/EthFollowup/Primes} (pure Mathlib; the files
compile with Lean v4.33.1 and v4.34.1) and \texttt{lean-bridge-anthropic/bridge/EthFollowup/Primes} (on top
of \A's development; the three pure files are linked into it). \Cref{tab:modules} lists the modules.

\begin{table}[ht]
\centering\small
\begin{tabular}{@{}p{0.40\textwidth}p{0.46\textwidth}r@{}}
\toprule
module & content & lines\\
\midrule
\lean{ChainParity} & chain lemma, one-count identity, $\pi(x)$ as a card & 249\\
\lean{ArithLDesc} & leveled descriptions and their count & 124\\
\lean{PrimeDesc} & the description \lean{primeDescL}, its count, well-formedness, spread & 589\\
\midrule
\lean{ArithLShape}, \lean{ArithLParams} & level shapes and boxes; wires, vertices, checks, canonical run & 1104\\
\lean{ArithLNetDefs}, \lean{ArithLNetParsimony} & the label system and its value & 643\\
\lean{ArithLNetDeg}, \lean{ArithLNetCount} & degrees; $U$ and $P$ (geometric sums) & 779\\
\lean{ArithLProcs}, \lean{ArithLPolyChk}, \lean{ArithLPoly} & procedures, checks on codes, polynomial time & 1425\\
\lean{ArithLPackage} & counting packages from leveled descriptions & 250\\
\lean{PrimePoly}, \lean{SpecPrime}, \lean{PrimePkgs} & readers, the one-query specification, the two packages & 279\\
\lean{Final} & the two final theorems & 46\\
\midrule
total & & 5488\\
\midrule
\lean{IterProdTree}, \lean{IterProdCutRank} & \cref{thm:tree}; rank of the product indicator (directory \texttt{Barriers}; not imported by \lean{Final}) & 142, 148\\
\lean{Check}, \lean{LemmaA} (in \texttt{Scratch/Primes}) & numerical checks by \lean{#eval}, not part of the proof & 66, 253\\
\bottomrule
\end{tabular}
\caption{The modules, with line counts (\texttt{wc -l}). The first three are pure Mathlib. The
development reuses, unchanged, the factoring modules of~\cite{TulFac26}, the ArithNet, package and
machine modules of~\cite{TulCN26} (\texttt{ArithDesc}, \texttt{ArithParams}, \texttt{GRate},
\texttt{ArithPkgs}, \texttt{Encode}, \dots) and the bridge modules of~\cite{Tul26}.}\label{tab:modules}
\end{table}

\paragraph{Trusted base and axioms.} \texttt{\#print axioms} reports \lean{propext},
\lean{Classical.choice} and \lean{Quot.sound} for the two theorems of \cref{thm:main}, for
\lean{ArithLSpec.pkg_exists}, and for the two theorems of \cref{thm:tree}. No module contains
\lean{sorry}, a new axiom or \lean{native_decide}.
A fresh \lean{leanchecker} replay of the dependency closure of \lean{EthFollowup.Primes.Final}, the
module holding the two final theorems, passed on 4~November 2026.
What must be trusted is the Lean kernel, the Mathlib definitions named in \cref{sec:statement}, our
definitions quoted there, and \A's definition of \texttt{Outputs}. \A's development, Mathlib's proofs and
our other modules are checked by the kernel; they are part of the checked closure, not of the trusted
base.

\paragraph{One design point.} Polynomial time must hold on \emph{every} context the reader may be given,
not only on the contexts it produces. The stored levels of a description are therefore clamped to the
number of variables before the cell layout uses them (\lean{ACtx.L} in \lean{ArithLProcs}); otherwise a
context with a large binary level would make the table procedure exponential.

% =====================================================================
\section{Related work}\label{sec:related}
% =====================================================================

We organise related work by idea: computing $\pi(x)$ exactly, parities on short intervals, deterministic
prime finding, and the counting engines with their formal proofs.

\paragraph{Computing $\pi(x)$ exactly.} The combinatorial method of Meissel and Lehmer, as extended by
Lagarias, Miller and Odlyzko~\cite{LMO85}, computes $\pi(x)$ in $O(x^{2/3+\varepsilon})$ steps, and
Del\'eglise and Rivat~\cite{DR96} save a factor $\log x$. The analytic method of Lagarias and
Odlyzko~\cite{LO87} takes $x^{1/2+o(1)}$ steps. Both give the parity as a by-product. In contrast, we
compute only the parity, but in $x^{o(1)}$ steps.

\paragraph{Parities on short intervals.} Tao, Croot and Helfgott~\cite{TCH12} prove that one can decide
whether an interval $[a,b]\subseteq[N,2N]$ of length at most $N^{1/2+c}$ contains an odd number of primes
in time $O(N^{1/2-c+o(1)})$ (their Theorem~1.2), from the identity
$\sum_{d^2\mid m}\mu(d)\tau(m/d^2)=2^{\omega(m)}$ and a fast divisor sum; they also treat the parity of
primes in progressions to small moduli. Lelechenko~\cite{Lel13} gives an explicit exponent
$O(x^{\max(c,7/15)+\varepsilon})$ for the same intervals (his Theorem~2). Sladkey~\cite{Sla12} computes the
divisor summatory function in $O(x^{1/3})$ steps and mentions the parity of $\pi$ among its applications,
citing a manuscript of Oliveira e Silva in preparation. Closest to our work is~\cite{TCH12}, whose
reduction of the parity to divisor counts we share. In contrast, we remove the M\"obius sum by
\cref{lem:chain} instead of evaluating it, which is what allows intervals of any length.

\paragraph{Deterministic prime finding.} Polymath4~\cite{Tao09} asked for a deterministic algorithm that
finds a $k$-bit prime in time polynomial in $k$, or at least $\exp(o(k))$; the best unconditional bound
remains $x^{1/2+o(1)}$, by computing $\pi$ analytically and bisecting~\cite{LO87,TCH12}. Scanning up from
$x$ with a primality test~\cite{AKS04} and the gap bound of~\cite{BHP01} gives $x^{0.525+o(1)}$. In
contrast, we give no prime finder, and \cref{sec:not} explains why the parity alone does not yield one.

\paragraph{Counting engines and formal proofs.} OpenAI~\cite{OpenAI26} and \A~\cite{Anthropic26} claim
deterministic subexponential $\#$SAT algorithms, with Lean developments, under review. Rudich and
Rousseau~\cite{RR26} factored in $N^{o(1)}$ with them, Shunia~\cite{Shu26} proved the rate $\rate$ on
paper, and~\cite{IF26} and our factoring note~\cite{TulFac26} proved it in Lean; our class-number
note~\cite{TulCN26} built the network library ArithNet, which we extend here by levels. We use
Lean~4~\cite{Lean4} and Mathlib~\cite{Mathlib20}. As of 4~November 2026 we found no derivation of a
parity result for $\pi$ from either ETH refutation. We searched the web for the parity of $\pi$ together
with ``ETH'', ``$\#$SAT'' and ``2026''; searched GitHub repositories and code for prime-counting parity
(the only hits were our own files); and checked the follow-ups known to us~\cite{RR26,Shu26,IF26}, which
treat factoring. Hexagon's listing is behind a login, so this negative is weak. In contrast to our earlier
notes, the number-theoretic content here is the chain lemma and the observation that working modulo
$2$ bypasses the M\"obius barrier; the rest is the network bookkeeping.

% =====================================================================
\section{Limitations}\label{sec:limitations}
% =====================================================================

\begin{itemize}
\item \emph{The result stands or falls with \A's development.} The proof rests on Anthropic's Lean
development of its claimed refutation of the Exponential Time Hypothesis, which is under review. Lean
checks that development as part of our proof, and our statements use only Mathlib, our quoted
definitions and \A's \texttt{Outputs}. We do need the development, built from source at the pinned
commit, to be sound, and the toolchain to be sound.
\item \emph{One bit only.} The machines output $\pi(x)\bmod2$ and $\pi([A,B])\bmod2$, not $\pi(x)$, and
not a prime. Residues modulo $k>2$ are on paper only (\cref{rem:modk}).
\item \emph{No prime finding.} The parity may be even on every interval one queries; \cref{prop:adv} is
an oracle-model statement on paper, and the barrier of \cref{sec:barrier} is proved only for product
trees.
\item \emph{The algorithms are galactic.} \A's levels start at $K_0=2^{x_0}$ with $x_0\ge18$, the level
needs $K^{40}\lesssim n$, and the constants ($2^{22}\cdot8$ for the package, and $C$) are astronomical. The
machines are of no practical use and beat $x^{1/2}$ only for astronomically large $x$.
\item \emph{The machines use exponential space.}
\item \emph{No human has checked the work.} The Lean code and this text were written with AI assistance.
No human expert has yet checked that \cref{sec:statement} lists everything a reader must trust.
\end{itemize}

\section{Conclusion}\label{sec:conclusion}
We have shown that one bit of $\pi$ on any interval is a single count of divisibility witnesses whose
sizes halve along a chain, and that the counting engine evaluates it in $x^{o(1)}$ steps once the network
pays for the witnesses at their true sizes. The result is a step towards, not a solution of, deterministic
prime finding. An interesting direction is whether the iterated product $m=d_1\cdots d_k$ with
$k=\Theta(n/\log n)$ has a parsimonious verifier with $o(n\sqrt{\log n})$ shared bits; a positive answer
would give the exact $\pi(x)$, and hence a prime in $[x,2x]$, in time $x^{o(1)}$, while a negative answer for
all networks would imply a circuit lower bound that is out of reach today.

\section*{Acknowledgements}
This work was done with AI assistance (Claude Code, Anthropic): the chain lemma, the Lean development,
the numerical checks, the literature and citation checks and the text were produced in sessions with
Claude Code. The Lean theorems are checked by the Lean kernel; everything else awaits human expert
review. Responsibility for any errors rests with the author. This note continues the project
of~\cite{Tul26,TulFac26,TulCN26}.

\begin{thebibliography}{99}

\bibitem{AKS04}
M.~Agrawal, N.~Kayal, and N.~Saxena, \emph{{PRIMES} is in {P}}, Ann. of Math. (2) \textbf{160} (2004),
no.~2, 781--793.

\bibitem{Anthropic26}
Claude (Anthropic), \emph{A $2^{o(n)}$ time algorithm for {3-SAT}}, preprint dated 26~October 2026.
Paper and Lean formalisation in the repository \url{https://github.com/anthropics/formal-math}, folder
\texttt{eth}; we use commit \texttt{16a5ce41304fcdaea4e0be7c132e1c361dc26da5}.

\bibitem{BHP01}
R.~C.~Baker, G.~Harman, and J.~Pintz, \emph{The difference between consecutive primes, {II}}, Proc.
London Math. Soc. (3) \textbf{83} (2001), no.~3, 532--562.

\bibitem{DR96}
M.~Del\'eglise and J.~Rivat, \emph{Computing $\pi(x)$: the {M}eissel, {L}ehmer, {L}agarias, {M}iller,
{O}dlyzko method}, Math. Comp. \textbf{65} (1996), no.~213, 235--245.

\bibitem{IF26}
1337isnot1337 (GitHub user) and contributors, \emph{integer-factoring: deterministic integer factoring
bounds}, public repository \url{https://github.com/1337isnot1337/integer-factoring}; Lean formalisation in
\texttt{formal-appendix/} (pull request \#20, merged 30~October 2026).

\bibitem{LMO85}
J.~C.~Lagarias, V.~S.~Miller, and A.~M.~Odlyzko, \emph{Computing $\pi(x)$: the {M}eissel--{L}ehmer
method}, Math. Comp. \textbf{44} (1985), no.~170, 537--560.

\bibitem{LO87}
J.~C.~Lagarias and A.~M.~Odlyzko, \emph{Computing $\pi(x)$: an analytic method}, J. Algorithms
\textbf{8} (1987), no.~2, 173--191.

\bibitem{Lean4}
L.~de Moura and S.~Ullrich, \emph{The {L}ean 4 theorem prover and programming language}, Automated
Deduction -- CADE 28, Lecture Notes in Comput. Sci., vol. 12699, Springer, 2021, pp.~625--635.

\bibitem{Lel13}
A.~V.~Lelechenko, \emph{Parity of the number of primes in a given interval and algorithms of the
sublinear summation}, Visn. Odes. Nats. Univ., Mat. Mekh. \textbf{18} (2013), no.~1(17), 104--114;
arXiv:1305.1639.

\bibitem{Mathlib20}
The mathlib Community, \emph{The {L}ean mathematical library}, Proc. 9th ACM SIGPLAN International
Conference on Certified Programs and Proofs (CPP 2020), 2020, pp.~367--381.

\bibitem{OpenAI26}
OpenAI, \emph{The {E}xponential {T}ime {H}ypothesis is false}, preprint dated 17~October 2026, revised
25~October 2026. Paper and Lean formalisation in the repository \url{https://github.com/openai/math}.

\bibitem{RR26}
I.~Rudich and L.-M.~Rousseau, \emph{Deterministic factoring in time $N^{o(1)}$}, Hexagon preprint
hexagon:2610.00600v1, published 29~October 2026, \url{https://hexagonmath.org/2610.00600v1}.

\bibitem{Shu26}
J.~Shunia, \emph{An improved subexponential bound for deterministic integer factorization}, Hexagon
preprint hexagon:2610.00711v1, published 2~November 2026, \url{https://hexagonmath.org/2610.00711v1}.

\bibitem{Sla12}
R.~Sladkey, \emph{A successive approximation algorithm for computing the divisor summatory function},
arXiv:1206.3369, June 2012.

\bibitem{Tao09}
T.~Tao, \emph{Deterministic way to find primes -- discussion thread}, The polymath blog, 28~July 2009,
\url{https://polymathprojects.org/2009/07/28/deterministic-way-to-find-primes-discussion-thread/}.

\bibitem{TCH12}
T.~Tao, E.~Croot III, and H.~Helfgott, \emph{Deterministic methods to find primes}, Math. Comp.
\textbf{81} (2012), no.~278, 1233--1246; arXiv:1009.3956 (as D.~H.~J.~Polymath).

\bibitem{Tul26}
V.~Tulsyan, \emph{Counting {H}amiltonian cycles in sparse graphs in subexponential time, assuming the
new $\#${SAT} algorithms}, Hexagon preprint hexagon:2610.00672v1, October 2026; source and Lean
development at \url{https://github.com/xennygrimmato/eth-hamiltonicity}.

\bibitem{TulCN26}
V.~Tulsyan, \emph{Class numbers, square roots and quadratic non-residues in time
$2^{O(n/\sqrt{\log n})}$, assuming the new $\#${SAT} algorithms}, preprint, November 2026; source in the
directory \texttt{classnumber-note/} of~\cite{TulsyanRepo26}.

\bibitem{TulFac26}
V.~Tulsyan, \emph{An independent machine-checked proof of deterministic factoring in time
$2^{O(n/\sqrt{\log n})}$, via the new $\#${SAT} algorithms}, preprint, November 2026; source in the
directory \texttt{factoring-note/} of~\cite{TulsyanRepo26}.

\bibitem{TulsyanRepo26}
V.~Tulsyan, \emph{eth-hamiltonicity: {L}ean 4 formalisation and research notes}, public repository
\url{https://github.com/xennygrimmato/eth-hamiltonicity}, 2026; Lean sources as of commit \texttt{f5aad1a}.

\end{thebibliography}

% =====================================================================
% CITATION CHECK (4 November 2026). Each entry was checked by web search or by fetching the source:
%  TCH12  arXiv 1009.3956v4 (26 May 2012), PDF read (pdftotext): Conjecture 1.1 (decide whether
%         [a,b] in [N,2N], length <= N^{1/2+c}, contains a prime in O(N^{1/2-c+o(1)})) open;
%         Theorem 1.2 "Computing the parity of pi(x)": odd number of primes in such an interval in
%         O(N^{1/2-c+o(1)}); sentence "we do not know of a method to quickly and deterministically
%         locate an interval with an odd number of primes" (sec. 1); Odlyzko N^{1/2+o(1)} in the
%         abstract; identity sum_{d^2|n} mu(d) tau(n/d^2) = 2^{omega(n)}; parity of pi(x;a,q) for
%         q = O(N^{c/10}). arXiv abstract page: journal ref Math. Comp. 81 (2012), 1233-1246, authors
%         listed in the journal as the main contributors (Tao, Croot, Helfgott). Issue no. 278 as in
%         the class-number note's audit.
%  Lel13  arXiv 1305.1639v3 (20 Sep 2013), PDF read: Theorem 2: [a,b] in [x,2x], b-a <= x^{1/2+c},
%         0 < c <= 1/2, parity of #{p in [a,b]} in O(x^{max(c,7/15)+eps}). Journal ref from arXiv:
%         Visn. Odessk. Univ., Ser. Mat. Mekh., 2013, vol. 18, #1 (17), 104-114.
%  Sla12  arXiv 1206.3369 (15 Jun 2012): O(n^{1/3}) time, O(log n) space; PDF read: ref. [5] =
%         "Tomas Oliveira e Silva, Efficient Computation of the Parity of the Prime Counting
%         Function, in preparation". Web search found no published version; we cite it only as
%         mentioned by Sladkey and give it no bibliography entry of its own.
%  LO87   J. Algorithms 8 (1987) 173-191 (reference list of TCH12; DOI 10.1016/0196-6774(87)90037-X,
%         June 1987); x^{1/2+eps} time per secondary sources (Platt, arXiv 1203.5712; DR96).
%  LMO85  Math. Comp. 44 (1985) no. 170, 537-560 (AMS page S0025-5718-1985-0777285-5);
%         O(x^{2/3+eps}) time, O(x^{1/3+eps}) space.
%  DR96   Math. Comp. 65 (1996) no. 213, 235-245, DOI 10.1090/S0025-5718-96-00674-6; abstract:
%         O(x^{2/3}/log^2 x) time.
%  BHP01  Proc. LMS (3) 83 (2001) no. 3, 532-562 (TCH12 reference list; Royal Holloway record);
%         [x - x^{0.525}, x] contains a prime for x > x_0.
%  AKS04  Ann. of Math. 160 (2004) no. 2, 781-793 (TCH12 reference list).
%  Tao09  polymathprojects.org post of 28 July 2009 by T. Tao (fetched): "Our goal is polynomial in
%         k"; "What about a partial result, such as exp(o(k))?". The michaelnielsen.org wiki page
%         "Finding primes" returned 404, so we cite the blog post.
%  Alboiu, Ramanujan J. 2015 ("On the parity of the prime-counting function and related problems",
%         DOI 10.1007/s11139-014-9596-1): about the distribution of pi(n) mod q, not computation;
%         not cited.
%  Anthropic26, OpenAI26, RR26, Shu26, IF26, Tul26, Lean4, Mathlib20: as audited in the factoring
%         and class-number notes (citation checks of 4 Nov 2026). TulFac26, TulCN26: our own notes,
%         unpublished.
%  TulsyanRepo26 commit: git log -1 --format=%h --
%         lean-bridge-anthropic/bridge/EthFollowup/Primes/Final.lean = f5aad1a.
%  Novelty searches (4 Nov 2026): web searches "parity of prime counting function 2026 #SAT ETH
%         refutation x^{o(1)} deterministic", "hexagonmath.org primes parity prime counting 2026",
%         "Oliveira e Silva parity of the prime counting function"; gh search repos "prime parity
%         ETH", "primeCounting parity lean"; gh search code "primeCounting x % 2" (only our own
%         files). Nothing relevant found.
% =====================================================================

\end{document}
